\section{Meditations on the Body}

\begin{quotex}
You were a model of perfection, full of wisdom, perfect in beauty; in Eden, the garden of God. \flright{\textsc{Ezekiel 23:12-13}}

\end{quotex}
I've been doing some preparation for the upcoming discussion on the sexes from an esoteric perspective. Besides the obvious sources, I started reading \textit{The Theology of the Body} by \textbf{Pope John Paul II}. I was more than a little surprised by what I read. As in any such work, the method is more important than the conclusion; the conclusion will make sense only by wrestling with the process.

Hence, this is a reflection just on chapter 1 which deals with the method. There are several themes which we have already addressed, showing the possibilities that can still be extracted from the Tradition. As such, the method is important in its own right, and not merely as a philosophically sophisticated justification of the doctrines on sexuality (the outcome is never in doubt). The real fruits will come when it is applied to other areas. I've highlighted several themes from my notes that warrant discussion.

\begin{itemize}
\item Opposition to modernity as defined by Bacon, Descartes, Kant 
\item Phenomenology as Hermetic meditation 
\item Subjectivity vs objectivity 
\item The relationship between experience and revelation 
\item Notion of the Person 
\item Nondual body/mind. Man as male and female 
\item Biblical texts as myth 
\item John of the Cross 
\item Self-mastery 
\item Tranquil witness of consciousness 
\end{itemize}
While a pope has to be circumspect about what he writes, we are free to be more speculative in drawing out the logical conclusions and place these ideas in a wider context.

\paragraph{Modernity}
The primary reason that a refutation of the modern project is impossible is that it is perfectly logical within its own presuppositions, much like an alternative geometry is to the Euclidean. If you change one axiom, the rest ineluctably follows. What is necessary, therefore, is the refutation of the axiom. For that task, the ratio, or rational mind is inadequate. Rather, only the intellectus, or higher mind can accomplish it. This is the ``intellectual conversion" that Guenon refers to. If the modern world is an illusion, only the truth can dispel it. But the truth must be ``seen" or ``intuited", not discovered by reason, or philosophy, or scientific experiment.

\textbf{Francis Bacon} inaugurated the modern project when he eliminated formal and final causes from his philosophy. As a methodological assumption, it turned out to be quite powerful, leading to scientific discoveries and technological advances. The problem arises when the Bacon's method is extended beyond that limited domain. Allegedly scientific theories about the origins of the universe, life, man, and society turn out to be destructive. The neglect of formal causes implies that the material world exists on its own and is self-explanatory. The denial of final causes implies that meaningless not just of the world, but also of man. There can be nothing ``natural" beyond what we experience because nothing has a purpose beyond that.

\textbf{Rene Descartes} took the next step. For him, the body is part of the material world and consciousness is detached from it, with no clear relationship. This leads to \textbf{Arthur Schopenhauer}'s famous question in the introduction to the World as Will and Idea: where exactly is the scientist in his theory? How does he explain himself? Specifically, where is he at the ``big bang" or the ``origin of the species"? He claims to be an objective observer at those events, but that is clearly an absurdity.

JPII challenges this project radically. The body, for him, is never inert matter, but rather a fully integrated body/spirit. The very experience of bodiness forms the basis for his meditations.

\paragraph{The Turn to Interiority}
Although JPII Does not reject the metaphysic of being, he is more interested in exploring the states of consciousness described in the Genesis story, rather than understanding it simply as objective fact. This has led to criticism in some quarters. However, the insistence on pure objectivity, if taken too far, implies that even atheists can be theologians (such is the case). Theology becomes a ``language game" with a set of rules that anyone can follow.

Now such an approach is not arbitrary subjectivity, from which many theological deviations arise. Rather, it is more like the way theology is understood in the East, as theoria, i.e., the theologian must actually ``know" God, not just talk ``about" God.

Thus, theology is not so much logic, but is based on the ``interior gaze" as JPII expresses it. Confusion may arise because the distinction between the psyche or soul and the intellectus or spirit is poorly understood. At the soul level, experience is indeed merely subjective. However, paradoxically, the spirit is totally objective. It is the ``witness to conscience" as JPII puts it.

The story of creation is ``prehistory", since it cannot be understood in terms of history or science. JPII describes it as a ``myth" and relies on figures such as \textbf{Mircea Eliade}, \textbf{Carl Jung}, among others to grasp mythological symbols in their fullness. Note that this does not mean JPII denies that something ``happened", but he proposes a way to understand it.

He calls this method phenomenology, but it is not your grandfather's phenomenology. Actually it is more like the Hermetic meditation described by \textbf{Valentin Tomberg}. He even refers to the lectures as ``meditations". So JPII uses the Biblical narrative to explore ever deeper levels of psychic layers. The main topics are ``original innocence", ``original solitude", the origin of shame, what it means to be embodied.

The objective text is insufficient. That is, revelation needs to be related to experience; it is not some incredible belief without relevance to life.

\paragraph{Sex}
Now Adam's original state was solitude. Although he named, or knew, all the animals, none were suitable companions to him. With Eve came the awareness of being embodied as male and female. Now this brings up an interesting point about how sex is related to humanity.

Definitions can be conventional or real. A conventional definition, for example, includes types of triangles (equilateral, isosceles). A type of triangle is still a triangle, it has the property ``triangularity". Hence, the distinction is conventional and is made for convenience and communication.

A real definition includes genus and difference, e.g., ``man is a rational animal". Man is of the genus ``animal", and what distinguishes him is ``rationality". So he is both animal and not-animal. Of course, the scientific project rejects that notion, and regards man as just a ``type" of animal. Rationality is a matter of degree not qualitative distinction.

So what then is sex? Is it more like a type of triangle? That is the modern view that regards one's sex as somewhat arbitrary, and a person can therefore reassign or redefine his own sexuality. That is rejected.

But is it not of the second kind, since both male and females are fully human, and not different species of the genus ``human"? JPII writes that sex is ``constitutive for the person" not an ``attribute of the person". Thus it is not conventional; it is essential to who one is, not accidental.

Thus there are differences between male and female. For example, man is the knower, woman is known, i.e., man is conscious.

\paragraph{Self-mastery}
In the Baconian project, every problem is presumed to have a technological solution. The body is no different from any other material object, so surgery, drugs, genetic manipulation, and other techniques are employed to solve existential issues.

However, for JPII, the human person is called to self-mastery. That is, the person is master of the body or, in Hermetic terms, the ``subtle rules the dense". Obviously this is also Evola's position when he writes of an inner autarchy.

Now to be free is to be free of constraints. JPII points out that a ``drive" is just an inner constraint. On the other hand, for the modern mind, freedom is the ability to act on one's ``instinctual" drives. Any ``voice" that opposes that misunderstanding is experienced as a threat. From that perspective, happiness should be the result of such acting out. In other words, ``the dense rules the subtle".

This false maxim is taken to the extreme by Marxism, for which the forms of consciousness are determined by the material aspects of life. Since few people bother to understand the root thoughts of their worldview, they may hold Marxist ideas even while denying it. For example, the notion that all social problems can be ``solved" by more education, programs, etc. is such a view. While a good regime is known for bringing peace and prosperity, it also requires a population that is able to exercise a certain measure of self-control.

\paragraph{Further Directions}
This method of phenomenology can be applied into new areas. Since the focus of \textit{The Theology of the Body} is on marriage and procreation, there is still much more to be said about being embodied as a man. There is the meditation on ``original innocence", as JPII says:

\begin{quotex}
Original innocence conceived in this way manifests itself as a tranquil witness of consciousness that precedes any experience of good and evil. 

\end{quotex}
Obviously, the ``witness" state is not descriptive, as phenomenology claims to be, but rather prescriptive. That is, it is a task and is therefore ``real" only for those able to achieve such a state. Elsewhere, he writes:

\begin{quotex}
Man is the causal origin of actions and at the same time the author of their meanings. 

\end{quotex}
Thus, a man must act consciously, understanding the meaning and consequences of what he does, for a determinate aim. Much of human activity is meaningless.

Similarly, rationality in the animal is virtual in most cases. A man's task, then, is to actualize that quality. We can quickly mention:

\begin{itemize}
\item A life lived in accordance with reason 
\item Becoming free from irrational drives, emotions, etc. 
\item Taking responsibility for one's actions 
\item Recognizing the personhood of others, not treating them as objects 
\item Achieving an inner state of calm and detachment 
\item Letting the subtle rule the dense 
\item Becoming the witness of consciousness 
\end{itemize}
These qualities, and more, deserve studies of their own.

\paragraph{I Sleep, but the Heart is Awake}

\begin{quotex}
I sleep, but my heart is awake. \flright{\textsc{Song of Solomon 5:2}}
\end{quotex}

Please consider the four states of consciousness mentioned by \textbf{Bede Griffiths} in \textit{The Cosmic Mystery}\footnote{\url{https://gornahoor.net/?p=8055}}.

We are accustomed to consider the sleeping, dreaming, and waking states as distinct states of consciousness, which follow each other consecutively. However, most of us have noticed that, rather than sequential, they interpenetrate each other. Our presumed waking state is mostly a series of semi-dreams: i.e., thoughts, imaginations, little stories, and so on. Then there is something deeper, beyond conscious awareness, in operation, affecting our postures, movements, attention and so on.

Among those dreams is the dream of awakening. That ``I", the subject of that dream, will ``snap us out of it" from time to time. That I needs to be nourished until the whole heart is awake.

\paragraph{Theology of the Body}
A delay was caused by some new information\footnote{\url{https://sspx.org/en/theology-of-the-body-genesis-of-confusion}} on the \textit{Theology of the Body} (H/T Perennial) that I came across as I was writing the introduction to the chapters on Sintesi that relate the sexes to race. Surprisingly, Evola's ideas have some things in common with \textbf{Pope John Paul II}'s conceptions. Perhaps that is so, given the influence of German idealism and the philosophy of personalism on both of them.

For Scholasticism, the mind, or intellectual soul, is the image of God. The mind is without qualities, neither male nor female. For JPII and Evola, on the other hand, the body is the outward expression of spirit, hence either masculinity or femininity permeates the whole being. The Scholastic approach easily devolves to a Cartesian dualism of body and mind.

Moreover, the Scholastic definition of the person is too abstract: \emph{an individual substance of a rational nature}. It does not mention Will, I, Freedom, Consciousness. Then there is Guenon's objection that it focuses on Being to the exclusion of non-Being. Specifically, it is too static, not making the distinction between the virtual and the actual. The actualization of a person is something else. Will must be developed, consciousness expanded, freedom/liberation must be achieved as it is not simply a given. ``Truth shall set you free." For the Scholasitics, there can be no development of the person. However, for JPII and Evola one of hallmarks of the person is self-possession, which requires a development.

Also, speculative theology does not make a clear distinction between soul and spirit. Mystical theology, on the other hand, does. \textbf{St John of the Cross}, for example, distinguishes between psychic experiences and truly spiritual experience. \textbf{St John of the Cross} is one of the influences on the Theology of the Body.

\paragraph{Sex and the Body}
These are some of the main points of the Theology of the Body. When the Evola translations come out, we can see better how they relate.

\begin{itemize}
\item There is no abstract ``human": the human being is always embodied as a man or a woman. Masculinity and Femininity are expressions of the person. 
\item The Spirit is not without qualities. 
\item The body is not just matter, subject to physical or biological laws. Rather, it is permeated by the soul and the spirit. Otherwise, resurrection cannot be understood except as the revival of a corpse. 
\item Biology is a human construct and does not exhaust all we know about the human being, sex, or the body. 
\item The Baconian project of gaining power over nature, specifically, technological power over the body, is rejected. Rather, the subtle rules the dense so the spirit must control the body, without relying on material means. 
\end{itemize}

\paragraph{Unintended Consequences}
There are perhaps other results that can be derived from this approach, not that they can be attributed to JPII. Biblical, mystical, philosophical, and phenomenological sources are used to develop the theology of the body. Why, then, can we not include other sources and even traditional sources? For example, Evola relies on \textbf{Otto Weininger}'s book \textit{Sex and Character} as well as the Laws of Manu for his own teachings on the spiritual and bodily aspects of masculinity and femininity.

Although JPII does not mention this, if the spirit has qualities, then race, or \emph{ethnos} as \textbf{E Michael Jones} calls it, has its counterpart in the spirit. The relationship is not necessarily one-to-one, as we will see. Of course, this doctrine rejects the idea that biology or DNA is the cause of differences in psychic and spiritual makeup of the person.

Furthermore, the body cannot be understood as a piece of matter, but rather as a non-dual body/soul/spirit complex. This conception is developed by \textbf{Valentin Tomberg} in Letter XX of the \textit{Meditations of the Tarot}. He shows the chain:

spirit $\Rightarrow $ psychic forces $\Rightarrow $ energy $\Rightarrow $ material organs.

That is the vertical process that works with the horizontal process of heredity. More research needs to be done to relate these various currents of thought.

Then, too, the notion that the body can be more or less dense begins to make some sense. Hence, there may be no physical traces of the Hyperboreans for the reason that the body in that era was less dense. This is something proposed by Rene Guenon. In this case, the positive sciences cannot be the last word.

\flrightit{Posted on 2015-05-08 by Cologero}

\begin{center}* * *\end{center}

\begin{footnotesize}\begin{sffamily}

\texttt{David on 2015-05-08 at 01:33 said: }

Fascinating to delve into the theology of recent popes. I also liked the extract you took from Benedict XVI some texts ago; he is a brilliant theologian. 

<<Now to be free is to be free of constraints. JPII points out that a ``drive" is just an inner constraint. Of course, for the modern mind, freedom is the ability to act out on one's ``instinctual" drives. Any ``voice" that opposes that misunderstanding is experienced as a threat. From that perspective, happiness should be the result of such acting out. In other words, ``the dense rules the subtle".>>

This also seems to be the error of psychoanalytic, i.e. the fascination with the subconscious as a force that drives and subject even the Self (considered in poor form, i.e. a social construct for overruling of subconscious urges). Psychoanalysis in more than one way is a direct inversion of Traditional principles, especially of Plato's ideas.


\hfill

\texttt{aegishjalmer on 2015-05-08 at 01:58 said: }

Thanks for that great analysis. Very helpful


\hfill

\texttt{Jacob on 2015-05-08 at 13:56 said: }

Good post Cologero. There need needs to be some clarity these days on how to think about sex because all other information is so biased and convoluted that it causes a lot of confusion. 

I struggle sometimes myself with eliminating Marxist materialist thoughts from my arguments. How can I oppose Liberalism as a force which causes the breakdown of Tradition without resorting to material conditions affecting a person's thoughts.? It goes back to a question I had earlier when Evola said that degeneration caused people to think a certain way and adopt certain ideologies. That seems almost like materialism to me. I had always assumed the opposite. Certain modes of thought was the cause of degeneration. Maybe there's a golden mean in between libertarian thought that the person is a self contained free willed individual with no external influences and Marxist determinism.


\hfill

\texttt{obscure on 2015-05-09 at 01:52 said: }

Cologero,

This was a very even post: A calm tide at the day's end reiterating what was also present before daybreak; only having achieved a greater height. The pregnant Moon lifts up the waters of the soul in the darkness and offers the Father's light and that very light pours in which is the Son and the image of the Father. The light travels far in the darkness but is never divided from the Father. The waters rise up in the darkness receiving the Son's mediation with the subtle aid of feminine gravity.

Then day breaks again and men forget the terrors and struggles of the night; falling down and away quickly until they are reminded again. The sea of souls rages in ignorance with the Father in full view. Although on some days there is a calm breeze of grace sent down from above. The light penetrating the breeze is the same Son carried in the womb of the Moon at high tide. And the breeze of air pulls the waters along the surface, carrying the fire and separating the salt; walking upon the waters and embodying the light of presence. What is far is not far and what is close is not close; what is close is far and what is far is close. The height creates the depth and the depth returns to the height because there is nowhere else to go upon the horizon. Anyone can see, if they want to, that those who chase the horizon never get to heaven. 

So, the wise stop moving their legs and instead they tilt back their heads. In this way, they understand that through not moving they are truly moving. This is why it is good to see good patterns in repetition.


\hfill

\texttt{Wayne Ferguson on 2015-05-09 at 15:10 said: }

John Paul II

``The Redemption of the Body and Sacramentality of Marriage" (Theology of the Body) 

``From the Weekly Audiences of His Holiness September 5, 1979 – November 28, 1984" 

\url{http://www.catholicprimer.org/papal/theology\_of\_the\_body.pdf}


\hfill

\texttt{X on 2015-04-22 at 01:16 said: }

``St John of the Cross, for example, distinguishes between psychic experiences and truly spiritual experience." 

Can you please say it more detailedly for this?


\hfill

\texttt{Cologero on 2015-04-22 at 05:36 said: }

X, this topic is discussed at length at \textit{Three Ages of the Interior Life}\footnote{\url{http://www.christianperfection.info/index.php}} by Fr Garrigou-Lagrange.

For example, concerning psychic phenomena: \textit{The Illuminative Way of Proficients}\footnote{\url{http://www.christianperfection.info/tta82.php}}.

Also, on mystical union, which can also be read as the antidote to Cassiodorus' objections: \textit{The Transforming Union}\footnote{\url{http://www.christianperfection.info/tta105.php\#bk2}}.


\hfill

\texttt{X on 2015-04-22 at 06:04 said: }

``The state which St. John of the Cross describes here is a state of love linked to a state of infused contemplation. The connection is owing to a necessity of love: `True and full love cannot hide anything.' This connection is not accidental, since this need is connatural to perfect charity."


\hfill

\texttt{thomas walker on 2015-04-22 at 10:11 said: }

We know from our own experience that body and soul are intimately related : bodily ailments can completely dominate our thoughts and emotions . I think that the kabbalistic idea of nephesh speaks to this , Our personalities are intimately linked with our bodies . There is a problem to understand how after death only the soul will exist until it is joined to the body in the Last Judgement . If time and space are abolished how can there be bodies which as bodies with extension must occupy space ?\newline
Also , doesn't Jesus teach that in heaven there will neither be taking or giving in marriage because bodies will be like the angels ( i assume that means without gender ) Therefore gender is only accidental not essential .


\hfill

\texttt{obscure on 2015-04-22 at 15:20 said: }

Nous is indeed prior to qualities, but qualities are received prior to the quantitative features of the body since quantity most signifies prime matter. The first relation to the body and thus to the macrocosm is qualitative, but pure Nous is not for this reason qualitative. Prime matter is also not seperable and thus not really distinct from information. Although form itself as Spirit or Idea is really distinct from matter or informed matter.

The rational soul is a discursive intellect animating a material substance or complex subject (these terms are synonymous). Its constitutive features are active intellect, passive intellect and sense-perception. Pure intellect is only an active intellect and its constitutive features are its simple form and its essential act (sometimes erronesouly translated by neo-Thomists in the tradition of Gilson as `act of existence'. but in truth `essential act' or `actus essendi' is pure will). A pure intellect is necessarily passive in relation to the First, but this is due to metaphysical necessity rather than discursive intellection with respect to the agent. The psychological (or `epistemological'. triad of the rational soul is not to be confused with the physiological triad (rational soul, sentient soul and vegetative soul) or the moral-constitutive triad (intellect, will and appetite).

Substantial form is the subsisting simple subjectivity of a complex subject (or compound substance, material substance, etc.: These terms are synonymous). A per se simple subject (not subsisting in a complex subject) is either a Spirit or Idea which is specifically different from the Absolute simple subject. Thus the metaphysical triad of all reality is Absolute simple subject, relative or `specific' simple subjects and lastly the complex subjects. In the order of pure simple subjects Spirits are active while Ideas are passive, although with respect to what is inferior (the order of complex subjects, the material cosmos, temporal procession, nature, etc.: These terms are synonymous) Ideas are obviously called `active'. but this is due to metaphysical necessity (which is to say, the hierarchy between God, supernature and nature). Corporeal form is constituted by extension and quality; it is an accidental form of a simple subject and it is the basic information of matter. Substantial form is not corporeal form, it `contains' essential possibilities not corporeal matter.

The basic qualification of Spirits other than the Absolute is their essential act or will gratuitously caused by the First (for the only relationship between the Infinite and the finite is gratuitous i.e. grace). The Absolute cannot be distinguished from His essential act in any respect since the Absolute is a Pure Act of Infinite Power. Even though neo-Thomists may not write in this way it is all to be found in the primary texts of their Doctor and of other medieval doctors. Duns Scotus is in many respects the master and I recommend this blog containing some historical materials if any are interested:

\url{http://lyfaber.blogspot.com/}


\hfill

\texttt{David on 2015-04-22 at 16:30 said: }

Cologero : Would it be possible to have all the texts within this site on a .zip or .rar format have a download link ? If the site is to end, a great part of the first texts I never had the chance to read. I imagine you have either a backup of .html files or the actual text files somewhere which could be given. Thank you.


\hfill

\texttt{obscure on 2015-04-22 at 18:48 said: }

Here are some translations of the Scotistic Theoremata:

\url{http://lyfaber.blogspot.com/2010/01/theoremata-scoti-partes-i-v.html}

Scotus' discusson of quality is interesting, although I'm sure it won't surprise many.


\hfill

\texttt{Cologero on 2015-04-22 at 20:28 said: }

TW, this is the full quote: ``In the resurrection they take neither wife nor husband, but are like the angels in heaven."

I suppose one approach is Guenon's: the being that was in the human state is now in an angelic state.

But. JPII contrasts the resurrection body with our bodies as they are now, with its members at war with the spirit. The resurrection body will not experience that opposition, it will be spiritualized. JPII is closer to Tomberg with this explanation:

\begin{quotex}
Spiritualization signifies not only that the spirit will master the body, but that it will also fully permeate the body and the powers of the spirit will permeate the energies of the body. (67:1) 

\end{quotex}
He takes pain to emphasize that this state is not ``disincarnation" of the body, nor a dehumanization. There is a deep harmony between spirit and body, but primacy resides in the spirit. He goes on to relate this to theosis (translated as divinization).

In the resurrection, the body is outside of human history ``tied to marriage and procreation". However, this does not mean that the body is suddenly neutered. (It would be interesting to compare that idea with the way Mouravieff includes marriage and procreation as part of the ``General Law".)


\hfill

\texttt{Cologero on 2015-04-22 at 20:47 said: }

Thanks for that synopsis, Obscure, but the issue is the whether or not JPII is introducing novelty into that doctrine. It appears that he is rejecting the notion that the Spirit is distinct from the body. For the Scotist, is masculinity or femininity an accidental or an essential quality? If God has qualities, then why would not a man have them?

As for your other link, I read the claim that the ``more perfect man has more perfect intelligence". I thought the Scotist position is that the goal of creation is the divinization of man, i.e., union with God. I don't see them as the same thing.


\hfill

\texttt{Cologero on 2015-04-23 at 08:55 said: }

David, Gornahoor will not be taken off-line, since it is intended for a future generation that may have a better perspective. However, I won't be updating it but perhaps someone else will.


\hfill

\texttt{obscure on 2015-04-23 at 23:15 said: }

Hello Cologero,

Keep in mind that when I write some long comment I intend a general audience and not just yourself in particular, for I enjoy sharing information. There is much ink I could spill with regard to the hierarchies of metaphysical qualification. This would require a fine treatment with some more effort than I can provide at the moment. If anything I've provided has produced some degree of confusion it was due to a degree of privation on my part. I promise to address your points at a later time.

As for the ToB, I can't say I've read it. My impression was that it was mostly an exoteric text of a more-or-less existential nature.

I will leave off for now with a Hermetic passage (Although it is not the sort of passage to be taken with little consideration or lack of interpretation):

``Soul enters the body by necessity, Intellect enters soul by judgment. While being outside of the body, soul has neither quality nor quantity. Once it is in the body it receives, as an accident, quality and quantity as well as good and evil for matter brings about such things."


\hfill

\texttt{Cologero on 2015-04-24 at 00:07 said: }

No confusion at all, Obscure, what you describe is how I would have been accustomed to look at it.

However, in light of the ongoing translation of Evola's Sintesi, there are questions that arise. I am not trying to provide the basis for a sophisticated theory of ``white nationalism", but in exploring the metaphysical ramifications of his doctrine. Precisely, is there spiritual differentiation among men?

It is not so obvious that corporeal life is an accident. For example, in Plato's Myth of Er, the circumstances of birth are determined by the prior state of the spirit. For Guenon, each human being is a possibility of manifestation in God's mind, and this determines the ``accidents" of birth. Perhaps these are differences in terms, but maybe not.

Based of the way it's been used, I, too, assumed that ToB was sort of a Catholic Kama Sutra. However, it is actually heavily indebted to John of the Cross. It also brings back final and formal cause to the understanding of the human being, contra both science and the entire modern project. It restores the notion of interiority to the body, which therefore cannot be treated simply as passive matter.


\hfill

\texttt{Michael on 2020-04-21 at 08:01 said: }

``Biology is a human construct and does not exhaust all we know about the human being, sex, or the body."

An interesting point when contrasted to the modern approach of ``constructs". While the modern approach says we can determine these for ourselves, JPII says there are higher methods of understanding to more great knowledge of the categories, a theology of the body (unified).

Yet in the other approach, we believe we can determine that human construct by just thinking it, isn't this just intellectual thinking centered from the individual? Pulling it down into discursive thought where all differences are forgotten instead of realized.

\hfill

\end{sffamily}\end{footnotesize}